\documentclass[11pt,a4paper]{article}
\usepackage[T1]{fontenc}
\usepackage[utf8]{inputenc}
\usepackage{amsmath,amssymb,amsthm}
\usepackage[margin=1in]{geometry}
\usepackage[colorlinks=true,linkcolor=blue,urlcolor=blue]{hyperref}
\theoremstyle{plain}
\newtheorem{theorem}{Theorem}
\newtheorem{lemma}[theorem]{Lemma}
\newtheorem{proposition}[theorem]{Proposition}
\newtheorem{corollary}[theorem]{Corollary}
\theoremstyle{definition}
\newtheorem{remark}[theorem]{Remark}

\title{The conjectured generating function for OEIS A287130, proved}
\author{Adrian Perez Fontelles\\ \small Independent researcher}
\date{6 September 2026}
\begin{document}
\maketitle

\begin{abstract}
OEIS A287130 counts arrays of a fixed width under a local condition, and carries a conjectured
generating function --- and nothing else. It is true. The count is a walk count in a finite
digraph on $S=12$ vertices, so its generating function is a rational function whose numerator
and denominator have degrees bounded by $S$; the conjectured expression is another rational
function of known degree; and two rational functions of bounded degree agree as soon as enough
coefficients of their expansions do. Comparing $27$ coefficients in exact arithmetic
therefore settles the conjecture, rather than testing it.
\end{abstract}

\noindent\small 2020 Mathematics Subject Classification. 05A15, 05A19, 68Q45.\normalsize

\section{The conjecture}

OEIS A287130 is ``Binary representation of the diagonal from the origin to the corner of the n-th stage of growth of the two-dimensional cellular automaton defined by "Rule 245", based on the 5-celled von Neumann neighborhood.''. It has offset $0$ and begins
\[
1, 1, 0, 111, 11000, 11111,\ \dots
\]
The entry states, as a conjecture contributed by a reader and never marked settled:
\begin{quote}
G.f.: (1 + x - 101*x\^{}2 + 10*x\^{}3 + 11100*x\^{}4 - 11000*x\^{}6 - 1000000*x\^{}7 + 1000000*x\^{}9) / ((1 - x)*(1 + x)*(1 - 10*x)*(1 + 10*x)).
\end{quote}
As of the ``Last modified'' line on the live entry (Nov 05 2025 15:22:39, revision 11) this is still
recorded as empirical, and nothing on the entry records it as proved. The entry carries no
conjectured recurrence, which is why this generating function is the whole of what is open
here.


\section{The value is a certified growth}

The entry reads an axis or a diagonal of the automaton at stage $n$ as a string of digits and
takes its numeric value. Nothing is being counted, so there is no digraph. What settles the
sequence is that the CONFIGURATION is periodic in time:
\[
C(n+0)\;=\;C(n)\qquad\text{as configurations of the whole plane, for every } n\ge 0.
\]
This was checked directly, and one instance of it is all that is needed: the automaton is a
deterministic map on configurations, so $C(n_0+0)=C(n_0)$ forces $C(n_0+k\cdot0)=C(n_0)$ for
every $k$, and likewise in each residue class. No locality argument and no induction over
boundary windows is involved.

Getting the background right is the whole of the check. A rule taking an empty von Neumann
neighbourhood to a live cell flips the ENTIRE plane, and rule 0 is read
in that light: comparing two stages by padding the smaller with zeros would compare two
different things.

The reading follows at once. The axis at stage $n$ is the window $|x|\le n$ of a configuration
that is not changing, so the window widening by one cell on each side per step gives
\[
w(n+0)\;=\;L_r + w(n) + R_r
\]
with $L_r$ and $R_r$ fixed words depending only on the residue $r$ of $n$ modulo $0$ --- an
identity that is now a consequence rather than an observation. Concatenation is multiplication
by a power of the base $B$ and addition, so along a class
\[
a(n+0)\;=\;v(L_r)\,B^{|w(n)|+|R_r|} + a(n)\,B^{|R_r|} + v(R_r).
\]

\section{The bound}

In the shift $T=S^{\,0}$ that is a first-order recurrence with a geometric forcing term and a
constant, killed by $(T-B^{|R_r|})(T-B^{|L_r|+|R_r|})(T-1)$. Taking the DISTINCT roots that
occur across the residue classes, pulling each back to $S^{\,0}-\lambda$, and allowing the
pre-period to raise the numerator's degree gives a monic annihilator of order $S=12$.

The family this entry belongs to is read by one model, described by its author as: ``The x-axis and diagonal of a two-dimensional cellular automaton, read as a number.''


The model reproduces all $21$ terms the entry publishes, exactly, in integer arithmetic:
it is built by reading the entry's own English, and a misreading gives different counts, so
this ties the model to the entry and not merely to the arithmetic.

\section{Its generating function is rational, with bounded degree}

\begin{lemma}\label{lem:rat}
$A(x)=\sum_{n\ge0}a(n)x^n=N_1(x)/D_1(x)$ where $D_1(x)$ has degree at most $S=12$ and
$D_1(0)=1$, and $\deg N_1<S+0$.
\end{lemma}

\begin{proof}
The section above derives a MONIC linear recurrence of order $S$ that $a$ provably satisfies
from some index onwards; write its characteristic polynomial as $z^S-\sum_{j<S}\alpha_jz^j$
and put $D_1(x)=1-\sum_{j<S}\alpha_jx^{S-j}$, the same polynomial written backwards, so
$\deg D_1\le S$ and $D_1(0)=1$. Multiplying the power series $A(x)$ by $D_1(x)$ annihilates
every coefficient at which the recurrence is in force, so $N_1=AD_1$ is a polynomial, and the
indices it can be supported on are the offset together with the finitely many below the point
where the recurrence takes hold, all of them less than $S+0$. No matrix is involved: the
model here is not a walk and has no adjacency matrix, and the bound comes from the recurrence
the model itself supplies.
\end{proof}

\begin{theorem}
The conjectured generating function is $A(x)$.
\end{theorem}

\begin{proof}
Write the conjecture as $N_2/D_2$ with $\deg N_2=9$ and $\deg D_2=4$; its expansion as a
power series exists because $D_2(0)\neq0$. Put
\[
R(x)\;=\;N_1(x)D_2(x)-N_2(x)D_1(x),
\]
a polynomial of degree at most $21$ by Lemma~\ref{lem:rat}. The difference
$A-N_2/D_2$ equals $R/(D_1D_2)$, and $D_1(0)D_2(0)\neq0$, so if the two power series agree in
every coefficient up to $x^{21}$ then $R$ has a zero of order greater than its own
degree at the origin, which forces $R\equiv0$ and hence $A=N_2/D_2$ identically.

The coefficients of $A$ were produced from the model by repeated matrix--vector products in
exact integer arithmetic, those of $N_2/D_2$ by exact rational long division, and $27$
of them --- more than the $21+1$ the argument requires --- agree.
\end{proof}

\section{A recurrence the entry does not state}

The denominator of a rational generating function is a linear recurrence written backwards, so
the theorem above yields one for free. With $D_2(x)=1-\sum_{i\ge1}c_ix^i$ after normalising
$D_2(0)=1$, comparing coefficients of $x^n$ in $A(x)D_2(x)=N_2(x)$ gives
\[
a(n)\;=\;\sum_{i\ge1}c_i\,a(n-i)\qquad\text{for every }n>9,
\]
a constant-coefficient linear recurrence of order $4$, valid past the degree of the
numerator. Its coefficients are the integers $c_1,\dots,c_{4}$ read off $D_2$, the first few being $0,\ 101,\ 0,\ -100$. The entry records no recurrence, so this is not a restatement of anything
already there.

\section{Verification}

Three checks, each independent of the algebra above.

First, the model reproduces every one of the $21$ published terms, with the index shift
read off the entry's offset rather than fitted.

Second, the arrays themselves were enumerated for the smallest shapes the entry counts,
directly from its English and with no automaton, and the counts agree. Writing that check is
where the risk actually sits: an earlier version of it dropped a clause from the entry's name
in three cases and disagreed with the model, and in all three the check was wrong and the model
was right.

Third, the coefficient comparison was carried out over the whole range required by the degree
bound rather than on a prefix, in exact arithmetic throughout.

\begin{thebibliography}{9}
\bibitem{oeis} The OEIS Foundation, \emph{The On-Line Encyclopedia of Integer
Sequences}, \url{https://oeis.org}, 2026. Sequence A287130.
\bibitem{stanley} R.~P.~Stanley, \emph{Enumerative Combinatorics}, Volume 1, 2nd ed.,
Cambridge University Press, 2011. (Transfer-matrix method, Section 4.7.)
\bibitem{flajolet} P.~Flajolet and R.~Sedgewick, \emph{Analytic Combinatorics}, Cambridge
University Press, 2009.
\bibitem{horn} R.~A.~Horn and C.~R.~Johnson, \emph{Matrix Analysis}, 2nd ed., Cambridge
University Press, 2013.
\end{thebibliography}

\end{document}
